% !TEX program = lualatex
% Transparencias de Fundamentos de las Matemáticas I
% Clase 04: Cuantificadores y negación

\documentclass[aspectratio=169,11pt]{beamer}

\usepackage{fontspec}
\usepackage[spanish,es-nodecimaldot]{babel}
\usepackage{amsmath,amssymb,mathtools}
\usepackage{booktabs}
\usepackage{csquotes}
\usepackage{xcolor}

\usetheme{Madrid}
\usecolortheme{default}
\setbeamertemplate{navigation symbols}{}
\setbeamertemplate{footline}[frame number]

\definecolor{FdMBlue}{RGB}{20,55,100}
\definecolor{FdMRed}{RGB}{130,30,30}
\definecolor{FdMGreen}{RGB}{25,100,60}

\setbeamercolor{structure}{fg=FdMBlue}
\setbeamercolor{title}{fg=white,bg=FdMBlue}
\setbeamercolor{frametitle}{fg=white,bg=FdMBlue}
\setbeamercolor{block title}{fg=white,bg=FdMBlue}
\setbeamercolor{block body}{bg=blue!3}
\setbeamercolor{alerted text}{fg=FdMRed}

\setmainfont{Latin Modern Roman}
\setsansfont{Latin Modern Sans}
\setmonofont{Latin Modern Mono}

\newcommand{\N}{\mathbb{N}}
\newcommand{\Z}{\mathbb{Z}}
\newcommand{\Q}{\mathbb{Q}}
\newcommand{\R}{\mathbb{R}}
\newcommand{\C}{\mathbb{C}}
\newcommand{\Pow}{\mathcal{P}}
\newcommand{\id}{\operatorname{id}}
\newcommand{\mcd}{\operatorname{mcd}}
\newcommand{\mcm}{\operatorname{mcm}}

\title[FdM I -- Clase 04]{Cuantificadores y negación}
\subtitle{Afirmaciones universales y existenciales}
\author{Antonio Falcó Montesinos}
\institute{Universidad CEU Cardenal Herrera}
\date{Fundamentos de las Matemáticas I}

\begin{document}

\begin{frame}
  \titlepage
\end{frame}

\begin{frame}{Objetivos de la clase}
\begin{itemize}
  \item Leer y escribir enunciados con \(\forall\) y \(\exists\).
  \item Negar correctamente enunciados cuantificados.
  \item Distinguir variables libres y ligadas.
\end{itemize}
\end{frame}

\begin{frame}{Mapa de la clase}
\tableofcontents
\end{frame}

\section{Cuantificadores}

\begin{frame}{Cuantificadores}
\begin{block}{Cuantificador universal}
\begin{itemize}
  \item \(\forall x\in A\, P(x)\): todos los elementos de \(A\) satisfacen \(P\).
  \item Para demostrarlo, se toma un elemento arbitrario de \(A\).
  \item Para refutarlo, basta un contraejemplo.
\end{itemize}
\end{block}
\end{frame}

\begin{frame}{Cuantificadores}
\begin{block}{Cuantificador existencial}
\begin{itemize}
  \item \(\exists x\in A\, P(x)\): hay al menos un elemento de \(A\) que satisface \(P\).
  \item Para demostrarlo, se construye o identifica un ejemplo.
  \item Para refutarlo, hay que probar que ningún elemento funciona.
\end{itemize}
\end{block}
\end{frame}

\begin{frame}{Cuantificadores}
\begin{block}{Orden de cuantificadores}
\begin{itemize}
  \item No es lo mismo \(\forall x\,\exists y\) que \(\exists y\,\forall x\).
  \item El orden cambia la dependencia entre objetos.
  \item En la lectura, conviene traducir primero a una frase completa.
\end{itemize}
\end{block}
\end{frame}

\section{Negación}

\begin{frame}{Negación}
\begin{block}{Reglas básicas}
\begin{itemize}
  \item La negación de ``para todo'' es ``existe un caso que no''.
  \item La negación de ``existe'' es ``para todo, no''.
  \item La negación atraviesa el cuantificador y cambia el cuantificador.
\end{itemize}
\end{block}
\end{frame}

\begin{frame}{Negación}
\begin{block}{Ejemplo}
\begin{itemize}
  \item Enunciado: \(\forall x\in\mathbb{R},\, x^2\geq 0\).
  \item Negación: \(\exists x\in\mathbb{R}\) tal que \(x^2<0\).
  \item No basta escribir ``no para todo'' sin precisar el caso contrario.
\end{itemize}
\end{block}
\end{frame}

\begin{frame}{Negación}
\begin{block}{Variables}
\begin{itemize}
  \item Una variable ligada está controlada por un cuantificador.
  \item Una variable libre necesita contexto externo.
  \item Antes de usar una fórmula, todas sus variables deben estar declaradas.
\end{itemize}
\end{block}
\end{frame}

\section{Profundización}

\begin{frame}{Profundización}
\begin{block}{Lectura por capas}
\begin{itemize}
  \item Primero se lee el cuantificador exterior.
  \item Después se identifica el conjunto al que pertenece la variable.
  \item Por último se lee la propiedad que se afirma de esa variable.
\end{itemize}
\end{block}
\end{frame}

\begin{frame}{Profundización}
\begin{block}{Cambio de orden}
\begin{itemize}
  \item \(\forall x\in\mathbb{R}\,\exists y\in\mathbb{R}: y>x\) es verdadero.
  \item \(\exists y\in\mathbb{R}\,\forall x\in\mathbb{R}: y>x\) es falso.
  \item El mismo vocabulario produce significados distintos al cambiar el orden.
\end{itemize}
\end{block}
\end{frame}

\begin{frame}{Profundización}
\begin{block}{Negación paso a paso}
\begin{itemize}
  \item Enunciado: para todo \(x\), existe \(y\) tal que \(P(x,y)\).
  \item Negación: existe \(x\) tal que para todo \(y\), no se cumple \(P(x,y)\).
  \item Cada cuantificador cambia y la propiedad final se niega.
\end{itemize}
\end{block}
\end{frame}

\begin{frame}{Cierre}
\begin{itemize}
  \item Negar cuantificadores es una habilidad central.
  \item El orden de los cuantificadores importa.
  \item La siguiente clase conecta lógica y métodos de demostración.
\end{itemize}
\end{frame}

\begin{frame}{Trabajo recomendado}
\begin{itemize}
  \item Releer las secciones correspondientes del manual.
  \item Identificar definiciones, símbolos nuevos y enunciados principales.
  \item Preparar dudas y ejemplos para el seminario asociado.
\end{itemize}
\end{frame}

\end{document}
